\section{Use cases}\label{sec:usecases}

\subsection{Overview and status}

Nine use cases: UC1--UC6 from 29 September, and UC7--UC9 for what was built since. Figure~\ref{fig:ucd} shows which persona takes part in which use case.

\begingroup\renewcommand{\arraystretch}{1.3}\small
\begin{longtable}{|>{\raggedright\arraybackslash\hspace{0pt}}p{0.9cm}|>{\raggedright\arraybackslash\hspace{0pt}}p{3.3cm}|>{\raggedright\arraybackslash\hspace{0pt}}p{2.6cm}|>{\raggedright\arraybackslash\hspace{0pt}}p{4.0cm}|>{\raggedright\arraybackslash\hspace{0pt}}p{3.0cm}|}
\hline
\hdr{ID} & \hdr{Use case} & \hdr{Primary actor} & \hdr{Goal} & \hdr{Status (30 Oct)} \\ \hline
\endfirsthead
\hline
\hdr{ID} & \hdr{Use case} & \hdr{Primary actor} & \hdr{Goal} & \hdr{Status (30 Oct)} \\ \hline
\endhead
UC1 & Ingest source data & Field Researcher, Lead Researcher & Bring the three perspectives' raw files into the system as usable input & \stBT \\ \hline
\rowcolor{light}
UC2 & Generate a draft themed section & Lead Researcher & Produce narrative text for one of the six themes, grounded in source data & \stBM \\ \hline
UC3 & Verify a claim against its source & Lead Researcher, Reviewing Colleague & Confirm or reject that a sentence is backed by the source it cites & \stB \\ \hline
\rowcolor{light}
UC4 & Correct or approve a claim & Reviewing Colleague & Edit, remove, or sign off a claim before the report is finalized & \stBT \\ \hline
UC5 & Configure model/provider & Lab Tech Lead & Choose which model runs each pipeline step and where data may go & \stP{} (config files; per run, not per step) \\ \hline
\rowcolor{light}
UC6 & Export the approved report & Lead Researcher & Produce a final report document matching InnoBeweegLab's report structure & \stBT \\ \hline
\addid{UC7} & Analyse the images of a case & Lead Researcher, Field Researcher & Turn map images and infographics into citable, checked evidence & \stBT{} (maps); vision reading: \stBM \\ \hline
\rowcolor{light}
\addid{UC8} & Record researcher notes & Field Researcher, Lead Researcher & Make team knowledge that no source file holds citable, with a source & \stBT \\ \hline
\addid{UC9} & Compare models against the original report & Lab Tech Lead & Measure which model and pipeline produce the best report & \stP{} (approaches E1 and E2; the checker variant E4 not built) \\ \hline
\end{longtable}
\endgroup

\begin{figure}[ht]
\centering
\begin{tikzpicture}[
  uc/.style={ellipse, draw=navy, line width=0.7pt, fill=white, text width=3.1cm, inner sep=0pt, minimum height=1.0cm, font=\footnotesize, align=center},
  ucadd/.style={uc, draw=amber, dashed, fill=amber!6},
  assoc/.style={draw=inkgray!70, line width=0.5pt},
  lblL/.style={font=\scriptsize, align=right, text width=1.45cm, anchor=east},
  lblR/.style={font=\scriptsize, align=left, text width=1.45cm, anchor=west}
]
\newcommand{\stick}[2]{%
  \draw[navy, line width=0.8pt] (#1,#2+0.45) circle (0.17);
  \draw[navy, line width=0.8pt] (#1,#2+0.28) -- (#1,#2-0.15);
  \draw[navy, line width=0.8pt] (#1-0.28,#2+0.12) -- (#1+0.28,#2+0.12);
  \draw[navy, line width=0.8pt] (#1,#2-0.15) -- (#1-0.22,#2-0.5);
  \draw[navy, line width=0.8pt] (#1,#2-0.15) -- (#1+0.22,#2-0.5);}
% system boundary
\draw[light!40!navy, line width=0.6pt, rounded corners=4pt, fill=light!40] (3.1,0.1) rectangle (13.0,9.0);
\node[font=\small\bfseries, text=navy, anchor=north west] at (3.2,8.95) {Reporting tool};
% use cases, left column
\node[uc]    (uc1) at (5.8,8.0) {UC1 Ingest source data};
\node[ucadd] (uc8) at (5.8,6.8) {UC8 Record researcher notes};
\node[ucadd] (uc7) at (5.8,5.6) {UC7 Analyse the images};
\node[uc]    (uc2) at (5.8,4.4) {UC2 Generate a draft section};
\node[uc]    (uc6) at (5.8,3.2) {UC6 Export the report};
\node[uc]    (uc3) at (5.8,2.0) {UC3 Verify a claim};
\node[uc]    (uc4) at (5.8,0.8) {UC4 Correct or approve};
% use cases, right column
\node[uc]    (uc5) at (10.5,6.2) {UC5 Configure model/provider};
\node[ucadd] (uc9) at (10.5,3.6) {UC9 Compare models against the original};
% actors (labels beside the figure, so no line crosses them)
\stick{2.1}{7.0}  \node[lblL] at (1.7,7.0) {Sam\\Field\\Researcher};
\stick{2.1}{4.4}  \node[lblL] at (1.7,4.4) {Anouk\\Lead\\Researcher};
\stick{2.1}{1.5}  \node[lblL] at (1.7,1.5) {Mika\\Reviewing\\Colleague};
\stick{13.7}{4.9} \node[lblR] at (14.1,4.9) {Alex\\Lab Tech\\Lead};
% associations
\draw[assoc] (2.4,7.12) -- (uc1.west);
\draw[assoc] (2.4,7.12) -- (uc8.west);
\draw[assoc] (2.4,7.12) -- (uc7.west);
\draw[assoc] (2.4,4.52) -- (uc1.west);
\draw[assoc] (2.4,4.52) -- (uc8.west);
\draw[assoc] (2.4,4.52) -- (uc7.west);
\draw[assoc] (2.4,4.52) -- (uc2.west);
\draw[assoc] (2.4,4.52) -- (uc6.west);
\draw[assoc] (2.4,4.52) -- (uc3.west);
\draw[assoc] (2.4,1.62) -- (uc3.west);
\draw[assoc] (2.4,1.62) -- (uc4.west);
\draw[assoc] (13.4,5.02) -- (uc5.east);
\draw[assoc] (13.4,5.02) -- (uc9.east);
\end{tikzpicture}
\caption{Use-case diagram. Dashed amber ellipses are the use cases added on 30 October (UC7--UC9). Lines connect each use case to the primary actors in the table above.}
\label{fig:ucd}
\end{figure}

% -----------------------------------------------------------------------------
\newcommand{\ucfield}[1]{\par\smallskip\noindent\textbf{\color{dknavy}#1}\quad}
\newcommand{\impl}[2]{\begin{calloutbox}{teal}{Implementation, 30 October: #1}#2\end{calloutbox}}

\subsection{UC1 --- Ingest source data}
\ucfield{Preconditions:} raw files exist in supported formats (coded/raw survey xlsx, expert score-and-remarks xlsx, QGIS map images, PDFs).
\ucfield{Main flow:}
\begin{enumerate}
\item Actor selects/uploads the source files for a given wijkscan.
\item System identifies file type and perspective (data/expert/bewoners).
\item System parses tabular files into rows tagged with theme, where a theme column already exists, or queues untagged rows for the tagging step.
\item System confirms what was ingested (counts per file/sheet) back to the actor.
\end{enumerate}
\ucfield{Alternate/exception flows:} an unsupported file format is rejected with a stated reason; a sheet whose structure matches no known template (neither ``Score per indicator'' nor ``Gecategoriseerde data'' shape, for example) is flagged for manual mapping rather than guessed at.
\ucfield{Postconditions:} all source rows/images are stored and individually addressable for later citation.

\impl{\stBT}{Oss (Wijkscan) and Buurtscan shapes: \code{ingestion/wijkscan.py}, \code{ingestion/buurtscan.py}, \code{ingestion/parsers/}. Every file gets a SHA-256 checksum, and byte-identical copies are skipped and reported (\code{ingestion/sources.py}). The command line and interface tab~1 show what was ingested per file. The original report PDF is registered as \code{excluded\_reference\_report}: it is used only for evaluation (UC9). Tests: \code{test\_buurtscan\_ingestion.py}, \code{tests/run\_stdlib.py}.}

\subsection{UC2 --- Generate a draft themed section}
\ucfield{Preconditions:} UC1 complete for at least one perspective touching the chosen theme.
\ucfield{Main flow:}
\begin{enumerate}
\item Actor picks a theme (e.g.\ ``Sporten'').
\item System retrieves the relevant rows/images across all ingested perspectives for that theme.
\item System drafts narrative text, attaching a source reference to each sentence/claim.
\item System marks any sentence it could not ground with a specific source.
\end{enumerate}
\ucfield{Alternate/exception flows:} no source data touches the chosen theme --- the system says so instead of drafting from nothing; the local model is unavailable or a hosted call fails --- the system reports the failure per section rather than silently omitting content.
\ucfield{Postconditions:} a draft section exists, unpublished, every claim carrying zero or more source links.

\impl{\stBM}{Step 2 runs through an evidence layer (\code{evidence/builder.py}): numbers are computed in code and each claim cites evidence ids, which lead to the cells. Pipelines \code{template}, \code{naive} and \code{staged} (\code{drafting/pipelines.py}, prompt version v2). Since 29 September the same step also drafts Conclusies, Aanleiding and Methode (D18--D20). Tests: \code{test\_staged\_runs\_one\_call\_per\_perspective\_plus\_synthesis}, \code{test\_empty\_pack\_does\_not\_draft\_from\_nothing}, \code{test\_invalid\_json\_is\_recorded\_not\_crashed}. Only the template has run on real data; the local models are not yet installed.}

\subsection{UC3 --- Verify a claim against its source}
\ucfield{Preconditions:} a draft section exists (UC2).
\ucfield{Main flow:}
\begin{enumerate}
\item Actor clicks a sentence/claim in the draft.
\item System shows the exact source excerpt(s) it was generated from --- a spreadsheet cell, a row's average score, or a highlighted region of a map image.
\item Actor judges whether the claim is supported.
\end{enumerate}
\ucfield{Alternate/exception flows:} a claim has no linked source (already flagged in UC2) --- the system shows ``no source found'' rather than an empty or misleading link.
\ucfield{Postconditions:} none required (a read-only check); often followed by UC4.

\impl{\stB}{Interface tab~3: per claim its check status, the reasons for any flag, and under ``Bewijs en bron'' each evidence item with its source rows (file, sheet, location, value, text). With no evidence it shows ``Geen bron gevonden''. Map evidence shows the measurement; the overlay image that shows what was counted is in tab~1b, not inside the review screen. No automated interface test; the underlying flag is tested (\code{test\_no\_citation}).}

\subsection{UC4 --- Correct or approve a claim}
\ucfield{Preconditions:} UC3 has been performed on the claim.
\ucfield{Main flow:}
\begin{enumerate}
\item Actor edits the claim's wording, deletes it, or approves it as-is.
\item System records the action, the actor, and a timestamp.
\item Once every claim in a section is approved or edited, the section is marked reviewed.
\end{enumerate}
\ucfield{Alternate/exception flows:} actor approves a claim the system had flagged as weakly grounded --- the system still records the override rather than silently accepting it, so the audit trail shows a human explicitly overrode the flag.
\ucfield{Postconditions:} the claim's status (approved/edited/removed) and the responsible actor are permanently recorded.

\impl{\stBT}{\code{review/session.py}: an append-only, hash-chained event log; the draft itself is never changed. Approving a flagged claim needs \code{confirm\_override} and an extra tick box in the interface. An edited claim is checked again. Regenerating a section archives the earlier review. Tests: \code{test\_flagged\_claim\_needs\_override}, \code{test\_edit\_keeps\_original\_and\_rechecks}, \code{test\_tampering\_is\_detected}, \code{test\_regenerate\_archives\_previous\_review}.}

\subsection{UC5 --- Configure model/provider}
\ucfield{Preconditions:} system installed on target hardware.
\ucfield{Main flow:}
\begin{enumerate}
\item Actor sets which model handles local tagging/extraction and which, if any, hosted API handles section drafting.
\item Actor sets what data classes may leave the device (e.g.\ anonymized text only).
\item System validates the configuration is internally consistent (e.g.\ refuses a config that would send unredacted resident data to a non-EU hosted API).
\end{enumerate}
\ucfield{Alternate/exception flows:} the chosen local model doesn't fit in available RAM --- the system warns before the actor commits the change.
\ucfield{Postconditions:} subsequent UC1/UC2 runs use the new configuration.

\impl{\stP}{Models are configured in \code{config/models.toml}: local or hosted, memory estimate, data residency. \code{llm/registry.py} refuses a local model above 60\% of 16~GB and a hosted provider without declared data residency, before a run. \code{llm/privacy.py} blocks resident free text to hosted providers unless an experiment sets a logged override, and pseudonymises first. Gaps: there is no settings screen, so configuration means editing a file (NFR10); the drafting model is chosen per run, not per pipeline step; and the only data class with its own rule is resident free text. Tests: \code{test\_model\_ram\_validation}, \code{test\_privacy\_blocks\_hosted\_resident\_text}, \code{test\_privacy\_override\_is\_logged}.}

\subsection{UC6 --- Export the approved report}
\ucfield{Preconditions:} all sections in scope are marked reviewed (UC4).
\ucfield{Main flow:}
\begin{enumerate}
\item Actor requests export.
\item System assembles the approved sections into the InnoBeweegLab report structure (Samenvatting $\rightarrow$ Methode $\rightarrow$ Resultaten per theme $\rightarrow$ Conclusies $\rightarrow$ Aanbevelingen $\rightarrow$ Bronnenlijst $\rightarrow$ Bijlagen).
\item System produces a document file.
\end{enumerate}
\ucfield{Alternate/exception flows:} a section is not yet fully reviewed --- the system blocks export of the full report and states which section is holding it up.
\ucfield{Postconditions:} a finished, human-approved report document exists outside the tool.

\impl{\stBT}{\code{export/report.py} and \code{export/docx\_report.py}: a Word report in the Wijkscan house layout (structure in \code{config/report\_structure.toml}), charts drawn from survey evidence, maps in the theme text, and an appendix that traces each claim to evidence and cell; also \code{evidence\_map.csv}. A preview for internal reading is marked CONCEPT. Recommendations are not generated; they stay with the researcher. Tests: \code{test\_export\_blocked\_until\_reviewed}, \code{test\_template\_run\_renders\_word\_report}.}

% ----------------------------------------------------------------------------- added
\subsection{UC7 --- Analyse the images of a case\added}
\ucfield{Primary actor:} Lead Researcher (confirms readings); Field Researcher (supplies the maps).
\ucfield{Preconditions:} UC1 has registered the case's map images and infographics; \code{config/images.toml} describes each image (kind, legend, themes, rules).
\ucfield{Main flow:}
\begin{enumerate}
\item Actor starts the image analysis for a case.
\item System measures each map in code: the share of coloured buildings per legend class and neighbourhood, the share of homes inside a distance zone, or the number of markers per type.
\item System writes an overlay image per measurement and stores each measurement as a citable record.
\item If a vision model is configured, the system reads the printed values on the infographics several times, keeps the median and runs range checks.
\item Actor confirms, corrects or rejects each reading; the system logs each decision in a hash-chained log.
\end{enumerate}
\ucfield{Alternate/exception flows:} a zone map cannot be aligned with the reference map (border agreement below 0.8) --- the measurement is refused, not guessed; two legend classes share a colour --- they are merged and flagged; no vision model is configured --- the infographics are listed as needing one, not skipped silently; the inspiration map shows proposals and is never measured as evidence.
\ucfield{Postconditions:} measurements and readings are evidence items that drafting can cite; a claim that cites an unconfirmed reading is flagged as weak; rejected readings are dropped.
\ucfield{Basis:} D11--D16, D26; Background Research v12, Sections 4.9 and 5.5.

\impl{\stBT{} for maps; \stBM{} for readings}{\code{imaging/maps.py}, \code{imaging/pipeline.py}, \code{imaging/vision.py}, \code{imaging/confirm.py}; interface tab~1b. Tests: \code{test\_imaging\_and\_oss.py}, \code{test\_notes\_zones\_conclusions.py}.}

\subsection{UC8 --- Record researcher notes\added}
\ucfield{Primary actor:} Field Researcher, Lead Researcher.
\ucfield{Preconditions:} a case exists (UC1).
\ucfield{Main flow:}
\begin{enumerate}
\item Actor creates the notes template (an Excel file) for the case.
\item Actor records each note with its kind (walk observation, resident meeting, stakeholder talk, project context, literature), its theme and its source.
\item System ingests the notes; each note gets an \code{O-\ldots} evidence id.
\item Drafting may cite notes; project context feeds the Aanleiding section and literature the Bronnenlijst.
\end{enumerate}
\ucfield{Alternate/exception flows:} re-ingesting the file replaces the notes instead of duplicating them; the template is never written over a filled file; notes from the resident meeting or a stakeholder talk carry the privacy flag and are attributed to the research team, never presented as the residents' view.
\ucfield{Postconditions:} team knowledge that no source file holds can be cited with its source.
\ucfield{Basis:} D17, D19; Background Research v12, Section 2.5 (in a sample of 54 statements from the reference report, 9 relied on such knowledge).

\impl{\stBT}{\code{ingestion/parsers/notes\_parser.py}, \code{note\_evidence} in \code{evidence/builder.py}; interface tab~1. Tests: \code{test\_template\_never\_overwrites\_filled\_notes}, \code{test\_ingest\_notes\_replace\_and\_survive\_reingest}, \code{test\_meeting\_note\_is\_not\_the\_residents\_view}.}

\subsection{UC9 --- Compare models against the original report\added}
\ucfield{Primary actor:} Lab Tech Lead, as researcher (SQ3).
\ucfield{Preconditions:} a case whose original human report is registered as reference and excluded from drafting; the local models are downloaded.
\ucfield{Main flow:}
\begin{enumerate}
\item Actor selects the models (\code{[selection]} in \code{config/models.toml}) and a pipeline.
\item System checks, before the run, that the model server runs and that each model is downloaded.
\item System drafts every section per model with the same evidence, prompts and pipeline, verifies it and renders a Word report; runs are kept apart from the review sessions.
\item System compares each report with the original: chapters, coverage of the original's statements, the ceiling the evidence can support, numbers found back, maps and charts; optionally with a local LLM judge.
\item System writes one comparison table, one row per run.
\item Actor labels a sample of statements; the system computes Cohen's $\kappa$ between the automatic and the human labels.
\item Actor may adopt one run as the starting point for the review (UC4).
\end{enumerate}
\ucfield{Alternate/exception flows:} a model cannot be reached --- the run is refused before it starts; a hosted judge is refused because the original report is confidential; the original's text never reaches a prompt.
\ucfield{Postconditions:} a comparison table per case; evaluation outputs that quote the original stay in the git-ignored workspace.
\ucfield{Basis:} D22--D25; SQ3; implementation brief Sections 15--17.

\impl{\stP}{\code{reports/runs.py}, \code{reports/evaluate.py}, \code{reports/judge.py}, \code{reports/cli.py}; interface tab~5. Tests: \code{test\_reports.py}. Not yet built: the checker pipeline (approach D), selection and seed options for full-report runs, averaging over repeats in the comparison table, and logging of the model's own memory use.}
